
\sourcenote{Genre: admissible rule (tier 2) --- a proof-tree-rewriting
recipe. The tier taxonomy defines tier 2 as \emph{verified unpacking
into core rules} (one node $\to$ many); this rule is the taxonomy's
recorded boundary case --- admissible by \emph{reduction to verified
machinery} --- and the recipe below typesets both the general template and
the honest degenerate instance.}

\begin{brecipe}[the admissible template, general form]\label{rec:pce-template}
An admissible rule $R$ with premises $P_1, \dots, P_n$ and conclusion $C$
is warranted by a \emph{rewrite}: a function taking any proof-tree node
\[
\inferrule*[Right=$R$]{\;\mathcal{T}_1 \rhd P_1 \\ \cdots \\ \mathcal{T}_n \rhd P_n\;}{C}
\qquad\rightsquigarrow\qquad
\mathcal{D}\bigl[\mathcal{T}_1, \dots, \mathcal{T}_n\bigr] \rhd C
\]
where $\mathcal{D}$ is a derivation fragment over core (and
previously-justified admissible) rules with the same conclusion, consuming
the premise subtrees $\mathcal{T}_i$ at its leaves. Applying the rewrite
leaf-first eliminates every $R$-node from a finite tree in finitely many
steps, so any judgement derivable with $R$ is derivable without it ---
admissibility --- and the calculus's soundness theorem is conserved with
\emph{no new kernel semantic content}. The verified object shipped per
admissible rule is either $\mathcal{D}$ itself or (equivalently, and how
this catalogue's Lean side does it) a soundness lemma for $R$ that factors
through already-verified machinery.
\end{brecipe}

\begin{brecipe}[the \rulename{PConstExact} instance --- the degenerate rewrite]
\label{rec:pce-instance}
\rulename{PConstExact} has \emph{no proof premises} ($n = 0$), and its
rewrite target contains \emph{zero catalogue nodes}: there is no core
derivation of a ground fact --- before this rule the calculus could not
state one at all ($\vdash -3/4 \leq 0$ had no derivation), which is the
rule's entire point. The rewrite therefore discharges the node against a
named \emph{verified theorem} rather than against catalogue structure:
\[
\inferrule*[Right=\rulename{PConstExact}]{\;\mathrm{evalClosed}(e_1 \bowtie e_2) = \mathsf{true}\;}{\Phi \vdash e_1 \bowtie e_2}
\qquad\rightsquigarrow\qquad
\text{\upshape the checker replays } \mathrm{evalClosed};\;
\text{\texttt{holds\_of\_evalClosed\_true} stands as warrant.}
\]
The side condition is \emph{computational} (exact $\Q$ evaluation of a
closed comparison --- decidable, replayed by the checker), and the warrant
theorem is kernel-verified once, globally: soundness of the rule is
literally ``the exact evaluator is correct.''
\end{brecipe}

\begin{bprop}[semantic warrant]\label{prop:pce-sound}
If $(\mathrm{atom}\ \bowtie\ e_1\ e_2).\mathrm{evalClosed} =
\mathsf{some\ true}$ and the sequent's conclusion is that atom, the
sequent is semantically valid over any hypernet, with arbitrary premise
$\Phi$.
\end{bprop}
\begin{proof}
Fix any valuation $\sigma$ (admissibility and $\sigma \models \Phi$ are
not consulted --- the rule is premise-free, the \rulename{Norm}
convention). By the verified evaluator theorem
\texttt{Formula.holds\_of\_evalClosed\_true}, a formula whose closed
evaluation returns $\mathsf{true}$ holds at \emph{every} valuation: a
closed term's evaluation consults no wire, so $\sigma$ is never read, and
the comparison was decided by exact arithmetic over $\Q$. That is the
whole proof --- \texttt{sound\_pConstExact} is one \texttt{exact} step.
\end{proof}

\begin{bremark}[closedness is implied, not declared]
There is no ``$e_1, e_2$ closed'' payload field to check:
$\mathrm{evalClosed}$ returns $\mathsf{none}$ on any wire occurrence, so
$\mathsf{some\ true}$ \emph{certifies} closedness as a byproduct. A
well-formedness hypothesis that the evaluator itself enforces is one the
payload cannot mis-declare --- unrepresentable, not discouraged.
\end{bremark}

\begin{bremark}[why $\neq$ is excluded --- the $M_3$ countermodel]
The comparator set is confined to $\{<, \leq, =, \geq, >\}$: five of the
six. The semantic lemma covers all six, but the calculus also maintains a
three-valued screen (\texttt{ProofObject.m3sound}), and under the $M_3$
countermodel a wire-free $\neq$ atom interprets to $\bot$ --- so admitting
$\neq$ against arbitrary $\Phi$ would let a premise-free rule conclude
something the $M_3$ screen refutes, breaking \texttt{m3sound} (the
\texttt{notNE} payload field, mirroring \rulename{ArithAtom}'s analogous
exclusion). A
hypothesis whose necessity comes not from $\Q$-semantics but from a
\emph{second} soundness invariant is exactly the kind a recipe must state,
or the next porter re-learns it the hard way.
\end{bremark}

\begin{bremark}[the subsumed corner, and coverage]
\rulename{PRatFact} (the scalar constant fact $c \bowtie 0$) is the corner
$e_2 = 0$ of this rule and was subsumed when the roster was consolidated. Coverage today:
the \texttt{Term} grammar makes every closed term rank-0, so the rule
states every closed \emph{scalar} exact comparison; the matrix corner
(exact closed matrix identities, MIX's use case) needs closed-tensor
constructors and a \texttt{QTensor}-valued closed evaluator --- the named
next leg, deliberately not hand-faked here.
\end{bremark}

\begin{anchors}
Recipe \ref{rec:pce-instance} warrant & \texttt{sound\_pConstExact} (one step to \texttt{Formula.holds\_of\_evalClosed\_true})\\
closedness-is-implied & \texttt{Formula.evalClosed} ($\mathsf{none}$ on any wire)\\
the $\neq$ exclusion & \texttt{notNE} payload field; \texttt{ProofObject.m3sound}\\
premise-freedom ($\Phi$ arbitrary) & the \rulename{Norm} convention; \texttt{pre} unused in \texttt{sound\_pConstExact}\\
wire placement & \texttt{RuleTag.pConstExact}; wire tag not yet allocated --- the emit path refuses the rule until allocation\\
\end{anchors}
